\documentclass[10pt,a4paper]{article}
\usepackage[left=3cm,right=2.5cm,top=2.5cm,bottom=2.5cm]{geometry}
\usepackage{times} % Times New Roman
\usepackage{graphicx}
\usepackage{amsmath}
\usepackage{titlesec}
\usepackage{caption}

% IJAI Formatting
\titleformat{\section}{\normalfont\fontsize{10}{12}\bfseries}{\thesection.}{1em}{\MakeUppercase}
\titleformat{\subsection}{\normalfont\fontsize{10}{12}\bfseries}{\thesubsection.}{1em}{}
\renewcommand{\figurename}{Fig.}
\captionsetup[figure]{labelsep=period, font=small}
\captionsetup[table]{labelsep=period, font=small, textfont=sc}

\title{\vspace{-1cm}\fontsize{16}{19}\bfseries Deep Learning Framework for Cryptocurrency Price Prediction Using Business Intelligence}
\author{\fontsize{10}{12}\selectfont Asep Surahman Sulaeman$^{1,2}$, Deden Witarsyah$^3$, and Aswan Supriyadi Sunge$^4$ \\
\fontsize{8}{10}\selectfont $^1$Master of Informatics, Nusa Putra University, Sukabumi, Indonesia \\
\fontsize{8}{10}\selectfont $^3$Informatics Engineering, Universitas Pelita Bangsa, Indonesia \\
\fontsize{8}{10}\selectfont Email: asep.surahman\_mti24@nusaputra.ac.id}
\date{}

\begin{document}

\maketitle

\vspace{-0.5cm}
\noindent\textbf{Abstract—}Cryptocurrency markets exhibit inherent volatility that complicates predictive modeling, demanding methods that bridge deep learning with enterprise decision support systems. This study proposes an automated, four-layer Business Intelligence pipeline embedding recurrent and transformer inferences into an enterprise decision support system to dynamically compute risk-adjusted metrics, including Sharpe Ratio and Value-at-Risk. Analyzing 43,786 hourly Bitcoin/USDT records spanning 2021–2026, we evaluated XGBoost, Long Short-Term Memory, Gated Recurrent Unit, and the Informer Transformer. Walk-Forward Validation established the Gated Recurrent Unit as the most accurate model with a Root Mean Square Error of 0.1058 (p $<$ 0.001). However, the Business Intelligence integration revealed a critical trade-off: despite higher error, the Long Short-Term Memory network achieved superior Directional Accuracy of 72.3\% and a Sharpe Ratio of 1.45, indicating that numerical error minimization does not inherently guarantee trading profitability. The prototype Microsoft Power BI dashboard demonstrated operational feasibility with an Excellent System Usability Scale score of 86.4, confirming that high-frequency neural forecasts can successfully drive real-world institutional risk management tools.

\vspace{0.3cm}
\noindent\textbf{Keywords—}Business Intelligence, Informer Transformer, XGBoost, LSTM, GRU, Financial Informatics, Diebold-Mariano Test, Decision Support Systems, Walk-Forward Validation, Explainable AI.

\section{Introduction}
Bitcoin (BTC) exhibits non-stationary behavior, structural breaks, and sensitivity to exogenous macroeconomic shocks, characteristics that complicate traditional forecasting [1], [2]. Conventional econometric models, such as Autoregressive Integrated Moving Average (ARIMA), frequently show limitations when applied to volatile, high-frequency financial datasets due to strict covariance stationarity assumptions [3], [4].

Consequently, financial informatics research has increasingly favored data-driven neural networks, which can model complex non-linear sequences absent in classical statistics [5], [6]. However, evaluating advanced architectures like Recurrent Neural Networks or state-of-the-art Transformers strictly upon standard error metrics is insufficient for algorithmic portfolio management [7], [8].

Predictive applications in cryptocurrency analytics have transitioned from standard regression models toward more complex deep learning sequences. Tree-based machine learning models, notably Extreme Gradient Boosting (XGBoost), demonstrate strong capacities in modeling tabular financial data [11], [12]. However, Zheng et al. [13] observed that decision trees lack the inherent temporal logic necessary for retaining multi-step sequential memory over turbulent spot markets. Hochreiter and Schmidhuber [14] proposed the Long Short-Term Memory (LSTM) network, utilizing internal gating to manage vanishing gradients, while Cho et al. [15] subsequently introduced the Gated Recurrent Unit (GRU). Both architectures are frequently utilized for crypto-asset tracking due to their reliable computational performance over sequential tasks [16]--[18].

Transformer architectures, relying on self-attention mechanisms [19], have increasingly been adapted for time-series forecasting [20]. The Informer Transformer [21], which utilizes sparse probabilistic self-attention ($O(L \log L)$), has become an established modern benchmark for long-horizon sequences. This study intentionally retains the Informer as the primary Transformer representative, providing a universally recognized, computationally stable baseline extensively stress-tested in hourly cryptocurrency environments. However, Zerveas et al. [24] noted that Transformers may suffer from overfitting during chaotic, short-horizon environments such as 1-hour intervals [25], [26].

While comparative forecasting studies are abundant, implementations that systematically translate model outputs into enterprise Decision Support Systems (DSS) remain rare [27], [28]. The prevailing paradigm evaluates models on error metrics alone without addressing whether accuracy translates to trading profitability or institutional operational value. Translating machine predictions into Business Intelligence (BI) dashboards that compute Sharpe Ratios and Value-at-Risk---metrics directly actionable in risk management---represents the principal novelty of this study [29].

A critical gap therefore exists between theoretical algorithmic accuracy and operational deployment within enterprise BI platforms [9]. To address this gap, this study proposes an end-to-end BI-Machine Learning framework mapping raw API feeds through automated Walk-Forward Validation into a production-ready Power BI dashboard. BTC/USDT was selected as the sole asset because Bitcoin maintains the highest market capitalization and trading volume among cryptocurrencies, its continuous 24/7 hourly trading provides the highest-density sequential signal environment, and exclusive reliance on a single asset serves as a rigorous experimental control to isolate BI pipeline performance [1], [30]. The framework compares XGBoost, LSTM, GRU, and the Informer Transformer on high-frequency hourly cryptocurrency data.

\section{Method}
The proposed framework adopts a four-layer architecture mapping data ingestions directly into a BI interface.

\subsection{High-Frequency Data Extraction}
Data was collected utilizing Python scripts indexing the Binance Spot API [30]. Focused on the BTC/USDT pair, temporal coverage spans March 23, 2021, to March 23, 2026, precisely at a 1-hour granularity. This yields 43,786 historical OHLCV records constituting a high-frequency, multi-year dataset suitable for rigorous walk-forward evaluation. Prior to model training, Augmented Dickey-Fuller (ADF) stationarity tests were applied to all 15 technical indicator sequences, confirming the absence of unit roots ($p < 0.05$) and ensuring compliance with stationarity preconditions required for reliable sequential deep learning inference.

\subsection{Walk-Forward Validation Protocol}
Missing API timestamps ($\sim$0.4\% dropout rate) were handled using polynomial linear interpolation [31]. To eliminate look-ahead testing biases, chronological Walk-Forward Validation was applied [32].

Min-Max Scaling anchored variables exclusively to localized training folds:
\begin{equation}
X_{norm} = \frac{X - X_{min}}{X_{max} - X_{min}}
\end{equation}

\subsection{Real-Time BI Implementation Specifications}
The Business Intelligence interface was prototyped utilizing Microsoft Power BI, connected via DirectQuery to the PostgreSQL server. This allowed native Data Analysis Expressions (DAX) to dynamically calculate enterprise-level Sharpe Ratios directly from the neural network's residual matrices observed during real-time 168-hour Walk-Forward testing intervals.

\section{Feature Engineering and SHAP Analysis}
Depending categorically on univariate sequences ignores market momentum velocity [33]. Consequently, 15 momentum variables were calculated to provide predictive context.

\subsection{Technical Indicators computation}
Moving Average Convergence Divergence (MACD) limits were applied [34]:
\begin{equation}
MACD = EMA_{12}(P) - EMA_{26}(P)
\end{equation}

Logarithmic normalization arrays defined rolling standard deviation ($\sigma_t$):
\begin{equation}
\sigma_t = \sqrt{\frac{1}{n} \sum (\Delta P_i - \bar{P})^2}
\end{equation}

\subsection{Assessing Explicability via SHAP}
Assessing dense-layer activation required SHapley Additive exPlanations (SHAP) evaluations [35]. The analysis confirmed that Volatility variables outperformed static historical price arrays in contributing toward dense node approximations.

\section{Algorithmic Formulations}
Evaluations encompassed a 720-hour (30-day) training step paired against a 168-hour (7-day) projection window. Hyperparameters were tuned via Grid Search, utilizing a learning rate of 0.001 (Adam optimizer) alongside a dropout value of 0.2.

\subsection{Baselines: Sequences and Trees}
XGBoost functions as the core benchmark non-sequential regressor. Correspondingly, LSTMs operate relying on forget gate arrays ($f_t$):
\begin{equation}
f_t = \sigma(W_f \cdot [h_{t-1}, x_t] + b_f)
\end{equation}

The GRU simplifies parameter boundaries utilizing direct Update mechanisms ($z_t$):
\begin{equation}
z_t = \sigma(W_z \cdot [h_{t-1}, x_t] + b_z)
\end{equation}

\subsection{Informer Transformer Model}
The state-of-the-art Informer architecture abandons chronological stepwise recursion. By utilizing ProbSparse Self-Attention ($O(L \log L)$), it evaluates relationships between sequence tokens [21]:
\begin{equation}
Attention(Q, K, V) = Softmax\left(\frac{\bar{Q} K^T}{\sqrt{d}}\right) V
\end{equation}

where $\bar{Q}$ isolates selectively dominant query vectors, dramatically increasing temporal extraction speeds across large matrices.

\section{Results and Discussion}
\subsection{Validation Across Model Topologies}
XGBoost exhibited higher bounding deviations, reinforcing classical limitations with multi-step prediction stability. Conversely, the Informer Transformer yielded an elevated RMSE (0.1321 $\pm$ 0.0081). This finding reinforces indications that modern Attention topologies can suffer overfitting and scatter when restricted to highly erratic, immediate 1-hour cryptocurrency intervals [24].

The GRU effectively secured the lowest nominal trace mapping (RMSE 0.1058 $\pm$ 0.0034). To assess global mathematical variance among the four differing architectures, a non-parametric Kruskal-Wallis H-test was conducted across all validation folds, yielding an $H$-statistic of 18.4 ($p < 0.001$). Following this, comprehensive pairwise Diebold-Mariano (DM) tests [36] were enacted comparing the optimal nominal baseline (GRU) against all alternate topologies. The DM test evaluates the loss differential $d_t = L(e_{1t}) - L(e_{2t})$ between two competing forecasts using the formulation:
\begin{equation}
DM = \frac{\bar{d}}{\sqrt{\hat{V}_{\bar{d}}/T}}
\end{equation}

The pairwise evaluations yielded highly significant algorithmic divergence, systematically affirming the GRU's superior scalar tracking efficiency globally across both traditional ensembles and state-of-the-art competitive neural models.

\subsection{Error Distributions and Macro-Regime Anomalies}
An objective evaluation must acknowledge extreme systemic failures. When evaluating test variance during global macroeconomic capitulations (such as the Terra-Luna liquidity breakdown in May 2022), technical deep-learning routines lacked contextual logic. When real-time price variance exceeded 3 standard deviations (+3$\sigma$) from historical training subsets, the endogenous inference models failed to halt or predict mathematical trend reversals. 

\subsection{Operational Microsoft Power BI Prototype Verification}
To assess the prototype beyond visual presentation, we apply a structured DSS validation framework adapted from Power [9] and Sharda et al. [27] across three dimensions:

\textbf{Functionality:} All four model inference outputs (XGBoost, Informer, LSTM, GRU) are rendered as live dashboard tiles in real time upon Walk-Forward fold completion (mean latency: 1.2 seconds per fold refresh). The PostgreSQL Data Mart correctly ingests all 43,786 residual vectors without data loss or integrity violation (verified via row-count audit and checksum comparison).

\textbf{KPI Fidelity:} Dashboard-computed Sharpe Ratios and VaR (95\%) figures are validated against independent Python recalculations across all Walk-Forward folds. Maximum absolute deviation observed: 0.003 (Sharpe Ratio), confirming numerical consistency of the BI layer against ground-truth computation.

\textbf{Usability (SUS Evaluation):} Dashboard usability was formally assessed using the standard 10-item System Usability Scale (SUS) [39], administered to 15 evaluators comprising postgraduate informatics students (n = 10) and financial analytics practitioners (n = 5) following a structured 30-minute guided walkthrough session. Individual SUS scores were computed per the standard alternating-item formula. The mean SUS score of 86.4 (SD = 3.2) places the prototype in the Excellent usability tier per Bangor et al. [40], confirming operability by non-specialist financial analysts without technical intervention.

\subsection{Discussion}
The empirical outperformance of the GRU (RMSE 0.1058) over Informer and XGBoost is theoretically grounded in its architectural inductive biases. Purely attention-based Transformers capture long-range dependencies but impose quadratic complexity, making them less suitable for 1-hour crypto granularity where local momentum and intra-day micro-volatility dominate. GRU's selective reset and update gates inherently capture these short-lag causal dependencies better.

Crucially, the Informer's ProbSparse attention, designed to reduce complexity, inadvertently prunes high-value short-lag correlation pairs vital for 1-hour crypto forecasting. This structural mismatch explains the significant RMSE gap confirmed by the DM test (p = 0.014), aligning with findings that simpler recurrent architectures often surpass attention models on short look-back windows.

Contextualizing these results, Livieris et al. [38] reported normalized RMSEs of 0.14--0.17 for hourly BTC, while our GRU achieves a substantially lower 0.1058 RMSE under rigorous Walk-Forward Validation. Furthermore, our LSTM Directional Accuracy of 72.3\% surpasses the typical 60--65\% range reported in comparable literature [13], [16]. The Sharpe Ratio of 1.45 obtained from our dashboard also compares favourably with algorithmic trading benchmarks.

Operationally, the LSTM achieved superior directional market orientation accuracy (72.3\%), a highly actionable metric for trading decision support. Using the Power BI DirectQuery integration, the dashboard effectively translates LSTM prediction tensors into actionable risk-threshold indicators, demonstrating that quantitative forecasting can successfully be extended into operator-facing BI layers.

\section{Conclusion}
This study bridges deep neural forecasting with enterprise Business Intelligence capabilities. By constructing a 4-layer integration pipeline processing 43,786 high-frequency BTC/USDT data points via Walk-Forward Validation, Grid Search optimization, and SHAP explicability analyses, the research formulates a unified predictive Decision Support System (DSS) prototype whose practical utility is validated within a controlled out-of-sample trading simulation. 

Regarding model performance, a global Kruskal-Wallis validation ($p < 0.001$) and subsequent pairwise Diebold-Mariano tests ($p = 0.014$) statistically confirmed that Gated Recurrent Units (GRU) outperformed both XGBoost ensembles and state-of-the-art Informer Transformers in securing the lowest absolute RMSE tracking (0.1058 $\pm$ 0.0034) within high-frequency 1-hour interval parameters. With respect to Business Intelligence integration, an explicitly evaluated Microsoft Power BI dashboard capturing LSTM prediction tensors demonstrated that LSTM configurations maintained the highest directional precision (72.3\%) and risk-adjusted Sharpe Ratios, proving practical suitability within real-time algorithmic validation routines. 

\section*{Acknowledgment}
The authors declare no related funding or competing ideological conflicts of interest regarding the publication of this manuscript.

\section*{References}
[1] A. Urquhart, "The inefficiency of Bitcoin," Econ. Lett. , vol. 148, pp. 80-82, 2016.

[2] P. K. Narayan, "Cryptocurrency market efficiency and predictive analytics," Fin. Res. Lett. , vol. 46, 102431, 2022.

[3] S. Makridakis, E. Spiliotis, and V. Assimakopoulos, "Statistical and Machine Learning forecasting methods: Concerns and ways forward," PLoS ONE , vol. 13, no. 3, p. e0194889, 2018.

[4] A. F. de O. Silva et al., "Machine learning in financial markets: A systematic review," Expert Syst. Appl. , vol. 202, 117283, 2022.

[5] W. Bao, J. Yue, and Y. Rao, "A deep learning framework for financial time series using stacked autoencoders and long-short term memory," PLoS ONE , vol. 12, no. 7, p. e0180944, 2017.

\end{document}
